% HD-F1-T36. GPT-6 Astra Pro, /root/transport; 22 November 2026 UTC.
\section[HD-F1-T36: bounded classical levels on the actual layer]{HD-F1-T36: a uniform effective level bound on every remaining large-rank symplectic or orthogonal constituent}
\label{sec:f1-classical-bounded-level-source-reduction}

Set
\[
                 J=10^{30},\qquad L_0=3J-3.
 \tag{F1T36.1}
\]
For every original remaining T2-OTHER source
with almost simple projective image and a
symplectic or orthogonal socle of rank
\(r\ge10^{31}\), every
irreducible constituent of the original
coefficient representation on the actual
natural classical layer has level
\[
                 \boxed{\quad 1\le\ell\le L_0.\quad}
 \tag{F1T36.2}
\]
The finite field still satisfies the retained
bound \(q\le27\). In odd orthogonal
characteristic T35 supplies the actual
natural quotient used here, including every
outer action. The bound (F1T36.2) is uniform
in \(q\) and in the original restriction
multiplicities. It leaves an unbounded natural
rank with a bounded level; it does not assert
that a remaining coefficient supplies a source.
We construct the precise bounded-tensor
inclusions at the end.

\paragraph{1. The actual classical layer and its outer actions.}
Retain \(G\subset GL_K(V)\),
\(m=\dim_KV\ge27\), its scalar centre \(Z\),
the inverse image \(P\) of the projective
socle \(S\), and \(L=P'\). Thus
\(P=ZL\), \(L=L'\), and \(Z(L)=L\cap Z\).
The universal-cover construction of T16,
T32 and T35 gives an actual surjection
\[
 \beta:H\longrightarrow L,\qquad
 H=\operatorname{Sp}_{2r}(q)
       \quad\hbox{or}\quad
 H=\Omega_N^{\epsilon}(q),
 \quad N=2r\text{ or }2r+1,
 \tag{F1T36.3}
\]
and, for every original \(g\in G\), an
automorphism \(a_g\) of this same \(H\) with
\[
                    \beta a_g=c_g\beta.
 \tag{F1T36.4}
\]
For the symplectic group, the simply connected
finite group is the displayed natural matrix
group. For an odd-characteristic orthogonal
group, (F1T36.3)--(F1T36.4) are exactly the
descended maps (F1T35.17)--(F1T35.18).
In even orthogonal characteristic the natural
central isogeny has kernel \(\mu_2\), whose
points over the algebraic closure consist of
the identity. Its natural finite image is
\(\Omega_N^{\epsilon}(q)\), and the induced
map on finite points is an isomorphism:
the unique lift of a Frobenius-fixed image
is Frobenius fixed. Thus every lifted
automorphism descends through that same
map. The centre and universal-cover
statements are the Malle--Testerman
Corollary~24.13/Table~24.2 and
Theorem~24.17/Remark~24.19 maps already
audited; all their exceptional ranks are
below the present range.

Here is also the exact matrix comparison for
an odd-dimensional orthogonal name in
characteristic two. Let \(A\) be that quadratic
space, \(R=\langle e\rangle\) the radical of
its polar form, and choose \(Q(e)=1\).
The quotient \(A/R\) is symplectic.
Every isometry fixes \(e\) and induces a
symplectic map on the quotient. This map of
isometry groups is injective: an element in
its kernel has the form \(v\mapsto v+f(v)e\),
and preservation of \(Q\) gives \(f(v)^2=0\).
It is surjective as well. Lift a symplectic
quotient map linearly, fixing \(e\). The
difference of the two quadratic forms has
zero polar form and vanishes on \(R\).
Over the finite perfect field it is the
square of a unique linear functional on
\(A/R\). Adding that functional times \(e\)
corrects the lift to an isometry. Uniqueness
makes these inverse constructions
homomorphisms. The quotient is therefore
the actual \(\operatorname{Sp}_{2r}(q)\)
matrix group; at this rank it is perfect,
so its full isometry preimage equals its
derived group. Conjugating every original
action through this isomorphism gives
(F1T36.4). No change of the original
coefficient representation is involved.

\paragraph{2. Actual isotypic projections force a common degree.}
Work on \(V_{\mathbf C}\) through the given
embedding \(K\hookrightarrow\mathbf C\).
For each irreducible \(L\)-character
\(\chi\) occurring in its restriction, use
the actual central idempotent
\[
 e_\chi=\frac{\chi(1)}{|L|}
       \sum_{x\in L}\chi(x^{-1})\rho(x).
 \tag{F1T36.5}
\]
The elementary character-orthogonality
relations make these orthogonal projections
whose sum on \(V_{\mathbf C}\) is the
identity. Conjugation by the actual matrix
\(\rho(g)\) carries \(e_\chi\) to
\(e_{\chi\circ c_g^{-1}}\). The sum of the
images in any one orbit is therefore an
actual \(G\)-submodule. Absolute
irreducibility forces a single orbit.
Consequently all occurring irreducibles
have the same degree \(d\), and conjugation
also preserves their multiplicities.

Inflation along the actual surjection
\(\beta\) preserves irreducibility and degree.
Writing the resulting \(H\)-modules as
\(W_i\), the canonical evaluation maps give
\[
 V_{\mathbf C}
    =\bigoplus_i W_i\otimes M_i,
 \quad M_i=\operatorname{Hom}_H(W_i,V_{\mathbf C}),
 \quad \dim W_i=d,
 \quad m=\sum_i d\,\dim M_i.
 \tag{F1T36.6}
\]
This is an actual decomposition for the
restriction, computed by (F1T36.5).
Equation (F1T36.4) specifies the conjugated
actions between its summands. A trivial
constituent would have a trivial conjugation
orbit and would force all of \(L\) to act
trivially, impossible for the literal
nontrivial subgroup \(L\subset G\).
Thus all \(W_i\) are nontrivial.
We shall use their common degree, not an
unstated assertion that an outer action
fixes a chosen Weil constituent or its level.

\paragraph{3. The exact published level inputs and their numerical consequence.}
Use Guralnick--Larsen--Tiep,
\emph{Inventiones mathematicae}
\textbf{235} (2024), 151--210,
Definition~3.2. Let \(A\) be the natural
space of dimension \(N\) for \(H\).
For \(H=\operatorname{Sp}_{2r}(q)\),
\(q\) odd, let \(\Theta\) be the actual
direct sum of the two reducible Weil
representations with characters
\(\omega_r,\omega_r^*\).
In the other cases let \(\Theta\) be the
direct sum with character \(\tau_A+\zeta_A\),
where \(\tau_A\) is the actual permutation
module on \(A\), and \(\zeta_A\) is the
specified restricted unitary Weil module
of that definition. The level \(\ell_i\)
is the least nonnegative integer for which
\(W_i\) occurs in \(\Theta^{\otimes\ell_i}\).
In particular it is the level on the actual
group (F1T36.3).

Lemma~3.4 of that paper gives
\[
 \begin{array}{c|c|c}
 H&N&\text{upper bound on }\ell_i\\ \hline
 \operatorname{Sp}_{2r}(q),\ q\text{ odd}&2r&2r+1\\
 \operatorname{Sp}_{2r}(q),\ q\text{ even}&2r&r+1\\
 \Omega_N^{\epsilon}(q)&N&\lfloor N/2\rfloor .
 \end{array}
 \tag{F1T36.7}
\]
Its Theorems~3.5--3.7 give the following
precise contrapositive at the fixed integer
\(k=J\). If \(\ell_i>3(J-1)\), then
\[
 d\ge
 \begin{cases}
 q^{rJ-J(J+1)/2}\bigl((q-1)/2\bigr)^J,
       &\operatorname{Sp}_{2r}(q),\ q\text{ odd},\\
 q^{2rJ-J(2J+1)}\bigl((q-1)^2/2\bigr)^J,
       &\operatorname{Sp}_{2r}(q),\ q\text{ even},\\
 q^{NJ-2J(J+1)}(q-1)^J,
       &\Omega_N^{\epsilon}(q).
 \end{cases}
 \tag{F1T36.8}
\]
The orthogonal exceptions \((N,k)=(8,2),(9,2)\)
do not occur. The full supplied proofs of
Lemma~3.4 and Theorems~3.5--3.7 were read,
including their actual parabolic restriction,
tensor and induction maps. These are named
published degree inputs, not the paper's
separate small-centralizer character estimates.

Suppose some actual constituent has
\(\ell_i> L_0\). Since the level is integral,
\(\ell_i\ge3J-2\). From (F1T36.7),
respectively,
\[
 r\ge(3J-3)/2\ge J+1,\qquad
 r\ge3J-3\ge2J+2,\qquad
 N\ge6J-4\ge4J+4.
 \tag{F1T36.9}
\]
All inequalities hold at the specified
\(J\), and the exponents below are positive.
In the first case \(q\ge3\) and
\((q-1)/2\ge1\), so
\[
             d\ge2^{rJ-J(J+1)/2}
                  \ge2^{rJ/2}=2^{NJ/4}.
 \tag{F1T36.10}
\]
In the second case \(q\ge2\) and
\((q-1)^2/2\ge1/2\), so
\[
             d\ge2^{2rJ-2J(J+1)}
                  \ge2^{rJ}=2^{NJ/2}.
 \tag{F1T36.11}
\]
In the orthogonal case \(q-1\ge1\), so
\[
             d\ge2^{NJ-2J(J+1)}
                  \ge2^{NJ/2}.
 \tag{F1T36.12}
\]
Thus in every case the common degree in
(F1T36.6) satisfies
\[
                         \frac{\log_2 d}{N}\ge\frac J4.
 \tag{F1T36.13}
\]
This conclusion applies to every actual
constituent, irrespective of any possible
permutation of their levels by the original
outer actions.

\paragraph{4. All-element ratios, original phases and the branch contradiction.}
Use the all-element Theorem~5.5 of
Larsen--Tiep, \emph{Annals of Mathematics}
\textbf{200} (2024), 1--70, with
\[
 \sigma=\frac1{6507\cdot2^{21}\cdot10^{18}}
                              >10^{-29}.
 \tag{F1T36.14}
\]
For a noncentral \(h\in H\), its natural
support is at least one: a scalar natural
matrix commutes with all of this faithful
matrix group. Hence, for every constituent,
\[
 \frac{|\chi_{W_i}(h)|}{d}
       \le d^{-\sigma/N}
       \le2^{-\sigma J/4}
       <2^{-5/2}<\frac14.
 \tag{F1T36.15}
\]
Summing the actual multiplicities gives
\(\left|\operatorname{Tr}_V\beta(h)\right|<m/4\).

For any original nonscalar \(g\in G\), the
trivial projective centralizer of the simple
socle supplies \(e\in H\) for which
\[
 h_g=a_g(e)e^{-1},\qquad
 \beta(h_g)=[g,\beta(e)]\notin Z(L).
 \tag{F1T36.16}
\]
Thus \(h_g\) is noncentral. Averaging a
Hermitian form over the original finite group
and using
\(\operatorname{Re}\operatorname{Tr}(u)
\ge m-2\operatorname{rank}(u-1)\)
gives
\(\operatorname{rank}(\beta(h_g)-1)>3m/8\).
The exact matrix commutator identity
(F1T35.14), with \(e_L=\beta(e)\), gives
\[
                       \operatorname{rank}(g-1)>3m/16.
 \tag{F1T36.17}
\]
Every nonidentity scalar has moving rank
\(m\). Thus every true original inertia
satisfies the weaker non-strict bound in
(F1T36.17), with no change of its phase.
The original product-one generating tuple
has \(b\ge12\) and total moving rank
\(2m+6\). Consequently
\[
                    2m+6\ge12\cdot\frac{3m}{16}
                         =\frac{9m}{4},
                    \qquad m\le24,
 \tag{F1T36.18}
\]
contradicting \(m\ge27\). Therefore no
actual constituent has level greater than
\(L_0\). Nontriviality in paragraph~2
excludes level zero and proves (F1T36.2).

\paragraph{5. Actual bounded-tensor inclusions and the retained source.}
For each constituent let
\(T_i=\Theta^{\otimes\ell_i}\).
Its isotypic projection is the literal
averaged operator
\[
 p_i=\frac d{|H|}\sum_{h\in H}
                    \chi_{W_i}(h^{-1})\rho_{T_i}(h).
 \tag{F1T36.19}
\]
The level definition says \(p_iT_i\ne0\).
The actual evaluation map
\[
 W_i\otimes\operatorname{Hom}_H(W_i,T_i)
                       \longrightarrow p_iT_i,
                    \qquad w\otimes f\longmapsto f(w)
 \tag{F1T36.20}
\]
is the semisimple isotypic isomorphism.
Choose any nonzero \(f_i\) in this Hom space.
Its kernel is an \(H\)-submodule of the
simple \(W_i\), so \(f_i\) is injective.
Together with (F1T36.6) this specifies an
actual \(H\)-equivariant inclusion
\[
 \iota:V_{\mathbf C}\longrightarrow
       \bigoplus_i T_i\otimes M_i,
 \qquad
 \iota|_{W_i\otimes M_i}=f_i\otimes1_{M_i},
 \quad 1\le\ell_i\le L_0.
 \tag{F1T36.21}
\]
On its image the original \(G\)-action is
the specified operator
\(\iota\rho(g)\iota^{-1}\). No action on
the unused tensor summands, or replacement
of the source coefficient by a whole tensor
power, is asserted. The map retains the
original multiplicity spaces and all
intertwiners between them.

The original cover, local system and inertia
tuple remain unchanged; their middle-extension
cohomology comparisons are identities.
The same is true of the original primitive
Jacobian corner, its integral-cleared image,
cup form and restricted polarization,
Poincar\'e biextension, Hodge filtration and
period differential. The Grothendieck--Adams
realizations of F01--F04 consequently have
the same identity comparisons and pairing
multiplier one. Equation (F1T36.21) records
the exact representation-theoretic bounded
level statement on that source. This strict
reduction contributes no detector component
and no unverified source existence statement.
